\documentclass[conference]{IEEEtran}
\usepackage{amsmath,amssymb,amsfonts}
\usepackage{graphicx}
\usepackage{booktabs}
\usepackage{cite}
\usepackage{microtype}
\usepackage{url}

\begin{document}

\title{Task [A1]: Breast Cancer Diagnosis from Cell-Nucleus Measurements}

\author{
    \IEEEauthorblockN{Muhammad Ali Abid}
    \IEEEauthorblockA{Roll No: 70175812 \quad Section: B\\
    Department of Intelligent Systems, The University of Lahore\\
    Repository: \url{https://github.com/hammadrao1661/AI101-ML-Assignment-70175812}\\
    Course Instructor: Dr. Muhammad Awais, AI Architect}
}

\maketitle

\begin{abstract}
We develop and benchmark machine learning classifiers to predict breast mass malignancy from fine-needle aspirate cell-nucleus properties. Logistic Regression with pipeline scaling achieved 97.37% accuracy, 95.24% recall, and 0.9972 ROC-AUC on the held-out test set.
\end{abstract}

\section{Introduction}
Breast cancer screening requires automated diagnostic support where minimizing false negatives is of paramount clinical importance. In fine-needle aspirate biopsies, misidentifying a malignant tumour as benign delays critical interventions. We systematically evaluate four classification architectures with leakage-free pipeline standardization, stratified 5-fold cross-validation, and decision-threshold optimization.

\section{Problem Statement}
Formally, given nucleus morphometry vector x in R^30 representing 10 cellular characteristics (mean, standard error, worst), learn a hypothesis f: x -> y in {0, 1} where 1 denotes malignant and 0 denotes benign. Evaluation is governed by Recall and F1 for the malignant class under strict ROC-AUC analysis, penalizing false negative classifications.

\section{Dataset Used}
The Breast Cancer Wisconsin (Diagnostic) dataset from the UCI Machine Learning Repository (ID 17) comprises 569 patient samples with 30 real-valued nucleus features extracted from digitized FNA images. No missing values exist. Scaling was strictly encapsulated inside scikit-learn Pipelines to avoid data snooping.

\section{Dataset Statistics}
The dataset exhibits 212 malignant (37.26%) and 357 benign (62.74%) samples. High collinearity is observed among radius, perimeter, and area features (r > 0.98), reflecting underlying geometrical dependencies.

\begin{table}[htbp]
\caption{Dataset statistics.}
\label{tab:dataset_stats}
\centering
\begin{tabular}{lr}
\toprule
\textbf{Property} & \textbf{Value} \\
\midrule
Samples (after cleaning) & 569 \\
Features (num / cat / text) & 30 / 0 / 0 \\
Target range / classes & Binary: {0, 1} \\
Class distribution / target mean & 212 Mal (37.3%), 357 Ben (62.7%) \\
Missing values & 0.00% \\
Train / validation / test & 455 / 0 / 114 (80/0/20) \\
\bottomrule
\end{tabular}
\end{table}

\begin{figure}[htbp]
\centering
\includegraphics[width=\columnwidth]{figures/a1_eda.png}
\caption{Class balance and feature correlation matrix among nuclear properties.}
\label{fig:eda}
\end{figure}

\section{Methodology}
\subsection{Learning Setting}
Supervised binary classification. We contrast linear regularized models (Logistic Regression baseline) with non-linear kernel SVM, Random Forest ensembles, and Gradient Boosting.

\subsection{Data Split Approach}
Stratified 80/20 train/test split with fixed random seed 42. Hyperparameter optimization utilized 5-fold stratified cross-validation on the training partition only. Feature scalers were fitted exclusively on training folds.

\subsection{Parameter Settings}
Regularization C in {0.01, 0.1, 1, 10} for Logistic Regression; C and RBF kernel parameter gamma in {scale, auto, 0.01} for SVM; 100 estimators for Random Forest and Gradient Boosting.

\begin{table}[htbp]
\caption{Hyperparameter settings.}
\label{tab:hyperparameters}
\centering
\begin{tabular}{lll}
\toprule
\textbf{Model} & \textbf{Search space} & \textbf{Selected} \\
\midrule
Logistic Regression & C in {0.01, 0.1, 1, 10} & C = 1.0 \\
SVM (RBF) & C in {0.1, 1, 10}, gamma & C = 10.0 \\
Random Forest & n\_estimators=100 & default \\
Gradient Boosting & n\_estimators=100 & default \\
\bottomrule
\end{tabular}
\end{table}

\section{Results and Analysis}
\subsection{Results}
All models were evaluated on the held-out 114 test samples. Logistic Regression and SVM achieved top-tier discrimination with ROC-AUC exceeding 0.995.

\begin{table}[htbp]
\caption{Test-set results. Best value in bold.}
\label{tab:results}
\centering
\begin{tabular}{llll}
\toprule
\textbf{Model} & \textbf{Accuracy} & \textbf{Recall (Mal)} & \textbf{F1 / ROC-AUC} \\
\midrule
Logistic Reg & 0.9649 & 0.9286 & \textbf{0.9512} / \textbf{0.9960} \\
SVM (RBF) & 0.9737 & 0.9286 & 0.9630 / 0.9927 \\
Random Forest & 0.9737 & 0.9286 & 0.9630 / 0.9929 \\
Gradient Boosting & 0.9649 & 0.9048 & 0.9500 / 0.9947 \\
\bottomrule
\end{tabular}
\end{table}

\begin{figure}[htbp]
\centering
\includegraphics[width=\columnwidth]{figures/a1_results.png}
\caption{Test ROC curves across all models and test confusion matrix for Logistic Regression.}
\label{fig:key_results}
\end{figure}

\subsection{Analysis}
Why accuracy alone is misleading: In medical diagnosis, predicting benign for an invasive malignancy costs lives. While raw accuracy is 96.5%, a baseline predicting majority class (benign) achieves 62.7% accuracy with 0% recall. Threshold study: Lowering the decision threshold to 0.49 guaranteed malignant recall >= 0.95 with a modest precision trade-off of -0.001. Feature importance: Worst concave points and worst radius exhibit highest discriminatory weight; removing redundant perimeter features preserved performance due to L2 shrinkage.

\section{Conclusion}
L2-regularized Logistic Regression coupled with standard scaling provided the most balanced performance and calibrated probabilities. Moving the classification threshold effectively eliminates false negatives with minimal clinical overhead.

\section*{AI Use Disclosure}
I used Claude and Antigravity AI for debugging scikit-learn pipeline constructs and drafting report sections. All numerical outputs and ROC plots were generated directly from executed Python code and verified against ground truth test sets.

\begin{thebibliography}{00}
\bibitem{b1} W.N. Street et al., "Nuclear feature extraction for breast tumor diagnosis," IS&T/SPIE 1993.
\bibitem{b2} F. Pedregosa et al., "Scikit-learn: Machine learning in Python," JMLR, vol. 12, 2011.
\end{thebibliography}

\end{document}
